\chapter{Varicose Veins}

\section*{Definition}
Varicose veins are dilated, tortuous, elongated superficial veins of the lower extremities, typically greater than 3 mm in diameter when measured in the upright position. Chronic venous disease is classified by the CEAP (Clinical, Etiologic, Anatomic, Pathophysiologic) system from C0 (no visible disease) to C6 (active venous ulcer).

\section*{What Happens in Your Body}
Venous hypertension from incompetent valves, venous wall weakness, or outflow blockage leads to venous reflux and pooling. Chronic venous hypertension causes capillary leakage of fibrinogen, hemosiderin deposition, and activation of inflammatory mediators that produce skin changes, swelling, and ulceration.

\section*{Causes}
\begin{itemize}
  \item \textbf{Primary (unknown cause):} Familial predisposition (most common), female sex, hormonal factors (pregnancy, estrogen), prolonged standing
  \item \textbf{Secondary:} Post-thrombotic syndrome (secondary to DVT), pelvic vein blockage, arteriovenous fistulas, obesity, advanced age
  \item \textbf{Occupational:} Occupations requiring prolonged standing (nurses, teachers, retail workers) or sitting (truck drivers, office workers)
  \item \textbf{Genetic:} Hereditary hemorrhagic telangiectasia, Klippel-Tr\'enaunay syndrome, Ehlers-Danlos syndrome, congenital valve absence
\end{itemize}

\section*{When to See a Doctor}
See your doctor if:
\begin{itemize}
  \item Progressive leg pain, aching, heaviness, or fatigue worsening with prolonged standing and relieved by elevation
  \item Skin changes: darkening of the skin, lipodermatosclerosis, stasis dermatitis, or atrophie blanche
  \item Venous ulceration (typically above the medial malleolus)
  \item Bleeding from a varicose vein (rupture)
  \item Suspected superficial thrombophlebitis or concomitant DVT
\end{itemize}

\section*{Related Symptoms}
\begin{itemize}
  \item Leg pain, heaviness, aching, fatigue, swelling, itching, restless legs, muscle cramps, burning sensation
\end{itemize}
\textbf{Important:} Varicose veins are a cosmetic concern for many, but the progressive nature of chronic venous insufficiency means symptoms typically worsen over time without intervention.

\section*{Self-Care / Home Management}
\begin{itemize}
  \item Compression stockings (20--30 mmHg or 30--40 mmHg grade) worn daily, putting them on before rising in the morning
  \item Leg elevation above heart level for 15--30 minutes, 3--4 times daily
  \item Regular exercise (walking, calf raises) to activate the calf muscle pump and promote venous return
  \item Weight loss to reduce intra-abdominal pressure and improve venous hemodynamics
  \item Avoid prolonged standing or sitting for more than 4 hours without breaks or movement
\end{itemize}

\section*{How Your Doctor Will Evaluate You}
\begin{itemize}
  \item Clinical examination in standing position: assess distribution of varicosities, swelling, skin changes, ulcers
  \item Duplex venous ultrasound (color flow) to evaluate deep and superficial vein patency, valve competence, and reflux duration
  \item CEAP classification: clinical class (C0--C6), causes (primary, secondary, congenital), anatomy (superficial, deep, perforator), pathophysiology (reflux, blockage, both)
  \item Venous clinical severity score (VCSS) to quantify disease severity and treatment response
  \item Venography or CT venography for pre-operative planning in complex cases or suspected pelvic source
\end{itemize}

\section*{Treatment Options}
\begin{itemize}
  \item Compression therapy (graduated stockings, multilayer bandaging) as first-line conservative management
  \item Endovenous ablation (laser, radiofrequency, or cyanoacrylate glue) for incompetent saphenous veins
  \item Sclerotherapy (liquid or foam) for reticular veins and telangiectasias
  \item Phlebectomy (ambulatory or stab avulsion) for large symptomatic varicosities
  \item Venous stenting for iliac vein blockage; wound care and compression for venous ulcers
\end{itemize}

\clearpage
